% =========================================================
% VII. DISCUSSION AND LIMITATIONS
% =========================================================
\section{Discussion and Limitations}
\label{sec:discussion_limitations}

The controlled experiment shows that structured event--rule retrieval and
selective verification improve exact decision labels. Full LegalSCM reaches
96.80\% accuracy versus 90.80\% for Causal Path RAG, with a statistically
significant paired gain. Since both systems expose identical retrieval and path
outputs, and all 17 Full-only gains occur after verifier commitment, the
difference is localized to committed factual validation and grounded
direct-edge checks.

The result does not imply a uniform quality gain. Answer F1 and ROUGE-L are
slightly lower for Full, and their paired differences, as well as the Citation
F1 difference, are not significant. Structured-label correctness and lexical
answer overlap therefore remain distinct objectives. Three erroneous cases
also contain prose that conflicts with the structured decision, motivating an
explicit generation-consistency constraint.

The closed-book LLM result is strongly prompt- and task-dependent. Its 0\%
accuracy arises because all 250 successful outputs abstain as
\textsc{uncertain}, while the gold set contains no abstention class. It is not
evidence that Qwen3-8B has zero legal knowledge. The 230/20 gold-label imbalance
also makes raw accuracy easy to inflate, so Balanced Accuracy, Macro-F1, class
coverage, and the confusion matrix should accompany it.

Most importantly, the reported experiment does not validate structural
mediator intervention. The 20 negative samples are direct-edge counterexamples,
and the saved Full run executes zero $do(X_M=0)$ operations. The 95\% subset
score measures grounded direct-edge topology. The hard-intervention mechanism
is implemented and formally separated from $G^{-M}$ node deletion, but requires
a new benchmark with explicit removal/necessity queries before effectiveness
claims can be made. Likewise, the engine supports conjunctive bodies and
alternative disjunctive mechanisms, whereas the current imported rules are
unary positive rules; rich AND/OR behavior is not empirically tested.

Top-1 Exact Path remains 38.56\% versus 75.42\% Oracle Exact Path. Seven of the
eight Full errors occur after a missing-primary-path fallback, confirming that
first-hop retrieval and path ranking remain bottlenecks. In addition, the
benchmark is small, graph-derived, and built from the same structured resource
used by the system. Manual expert review, out-of-graph legal questions, and
external-domain testing are needed for external validity.

Finally, LegalSCM depends on rule extraction, event normalization, endpoint and
mediator grounding, and graph construction. It is a deterministic normative
rule model, not causal discovery or effect estimation. Its $do(\cdot)$ operator
is meaningful only relative to encoded mechanisms and lacks the exogenous
abduction layer required for full unit-level counterfactual analysis.

% =========================================================
% VIII. CONCLUSION
% =========================================================
\section{Conclusion}
\label{sec:conclusion}

LegalSCM-RAG combines causal-path retrieval, three-valued rule inference, and a
selective commit gate for Vietnamese multi-hop legal QA. On the current
250-question benchmark, it improves decision accuracy over Causal Path RAG by
6.0 points with paired statistical support, but does not significantly improve
answer or citation metrics.

The strongest supported conclusion is therefore that selective factual and
direct-claim verification improves decision reliability in this controlled
setting. Future work will construct explicit mediator-intervention and
multi-antecedent rule benchmarks, improve first-hop retrieval and endpoint
grounding, enforce answer--decision consistency, and obtain expert evaluation
on legal questions not generated from the source graph.
